\documentclass[11pt]{article}
\usepackage[margin=1in]{geometry}
\usepackage{amsmath,amssymb,amsthm,mathtools}
\usepackage{bm}
\usepackage{booktabs}
\usepackage{graphicx}
\usepackage{enumitem}
\usepackage[colorlinks=true,linkcolor=black,citecolor=black,urlcolor=blue]{hyperref}

\newcommand{\R}{\mathbb{R}}
\newcommand{\Gr}{\operatorname{Gr}}
\newcommand{\RP}{\mathbb{RP}}
\newcommand{\Fld}{\mathbb{F}}
\newcommand{\tr}{\operatorname{tr}}
\newcommand{\rank}{\operatorname{rank}}
\newcommand{\Log}{\operatorname{Log}}
\newcommand{\Exp}{\operatorname{Exp}}
\newcommand{\PT}{\mathcal{P}}
\newcommand{\norm}[1]{\left\lVert #1 \right\rVert}
\newcommand{\inner}[2]{\left\langle #1,\, #2 \right\rangle}
\newcommand{\Span}{\operatorname{span}}
\newcommand{\diag}{\operatorname{diag}}

\theoremstyle{definition}
\newtheorem{definition}{Definition}
\newtheorem{remark}{Remark}

\title{\textbf{Subspace Geometry and Topology of High-Dimensional Neural Time Series}\\[4pt]
\large A Grassmannian pipeline: definitions, constructions, and estimators}
\author{Projective Space Models --- methods note}
\date{\today}

\begin{document}
\maketitle

\begin{abstract}
\noindent
We describe a pipeline that characterises the \emph{geometry} and \emph{topology} of
high-dimensional neural time series (e.g.\ recurrent-network activity or population
recordings) by tracking, over time, the low-dimensional linear subspace that the
activity locally occupies. Each local subspace is realised as a point on the
Grassmannian manifold $\Gr(k,N)$; a time series therefore becomes a trajectory on
$\Gr(k,N)$. We then quantify (i) the Riemannian \emph{kinematics} of that trajectory
(speed, covariant acceleration, curvature), (ii) intrinsic \emph{statistics} of the
subspace distribution (Fr\'echet mean, tangent-space PCA), (iii) the intrinsic
\emph{differential structure} of the sampled submanifold via discrete exterior
calculus (DEC), and (iv) the \emph{topology} of the trajectory and of the pooled
point cloud via persistent homology over $\Fld_2$. For $k=1$, $\Gr(1,N)$ is real
projective space $\RP^{N-1}$, which motivates the name of the framework. This note
fixes notation, states every construction as a definition, and derives the
estimators used downstream. A companion critique note discusses assumptions and
limitations; a companion slide deck reports the empirical results on a synthetic
recurrent-network testbed.
\end{abstract}

\tableofcontents

\section{Overview and motivation}
\label{sec:overview}

Let $\bm{x}(t)\in\R^N$ denote the activity of $N$ units at time $t$. The standard
state-space view studies the raw trajectory $t\mapsto\bm{x}(t)$ directly, e.g.\ via
PCA, and asks about the shape of the swept manifold. The premise here is
complementary: much of the computation performed by a population is expressed not
by \emph{where} the state is, but by \emph{which} low-dimensional subspace the
activity currently occupies, and by \emph{how that subspace moves}. Concretely,
over a short window the activity fluctuates within a $k$-dimensional linear subspace
$\mathcal{S}(t)\subseteq\R^N$; the computation is read out from the temporal path
$t\mapsto\mathcal{S}(t)$ in the space of all $k$-dimensional subspaces.

The set of $k$-dimensional linear subspaces of $\R^N$ is the Grassmannian
$\Gr(k,N)$, a compact Riemannian manifold of dimension $k(N-k)$. Mapping activity to
$\Gr(k,N)$ turns a neural time series into a curve on a curved manifold, so that (a)
the \emph{rate and manner} of subspace reconfiguration become intrinsic Riemannian
quantities, and (b) recurrence and periodicity of the computation become
\emph{topological invariants} (loops, voids) of the curve or of the pooled ensemble
of subspaces. Two lenses are cross-checked throughout: the \emph{subspace lens}
(activity $\to$ Grassmannian trajectory $\to$ kinematics/DEC/TDA) and a \emph{direct
lens} (activity $\to$ state-space manifold estimate $\to$ TDA), which guards against
metric artifacts.

\section{Data model and notation}
\label{sec:data}

\paragraph{Generative model (testbed).} Synthetic data are generated by a
continuous-time rate RNN,
\begin{equation}
\tau\,\dot{\bm{x}}(t) = -\bm{x}(t) + W\,\phi\!\left(\bm{x}(t)\right) + B\,\bm{u}(t) + \bm{\xi}(t),
\qquad \phi=\tanh,
\label{eq:rnn}
\end{equation}
with time constant $\tau$, recurrent weights $W\in\R^{N\times N}$, input $\bm{u}(t)$
entering through $B$, and (optional) white-noise current $\bm{\xi}$. Deterministic
runs use fixed-step RK4; noisy runs use Euler--Maruyama, so that a step of size $dt$
adds $\sigma\sqrt{dt}\,\bm{\eta}_k$ with $\bm{\eta}_k\sim\mathcal{N}(0,I)$. The
pipeline is agnostic to \eqref{eq:rnn}: any array of states in the schema below is
admissible, including task-trained networks and in-vivo recordings.

\paragraph{Data tensor.} A dataset is $X\in\R^{n_{\mathrm{tr}}\times T\times N}$
($n_{\mathrm{tr}}$ trials, $T$ time samples, $N$ units) with sample times
$\{t_1,\dots,t_T\}$. A single trial is $X^{(r)}\in\R^{T\times N}$, whose row
$\bm{x}^{(r)}(t_i)$ is the state at sample $i$.

\paragraph{Linear algebra.} For $A\in\R^{N\times k}$ with orthonormal columns
($A^\top A = I_k$) we write $\Span(A)$ for its column space. $\norm{\cdot}_F$ is the
Frobenius norm. The singular values of a matrix $M$ are $\sigma_1(M)\ge\sigma_2(M)\ge\cdots$.

\section{Sliding-window subspace embedding}
\label{sec:embed}

We convert a single-trial trajectory into a sequence of subspaces by a sliding-window
singular value decomposition (SVD).

\begin{definition}[Window frame]
\label{def:frame}
Fix a subspace dimension $k$, window length $w$ (samples) and stride $s$. For window
$m$ starting at sample index $a_m = 1 + (m-1)s$, form the block
$X_m \in \R^{w\times N}$ of the $w$ consecutive state rows. Optionally centre it in
time, $\tilde{X}_m = X_m - \mathbf{1}\bar{\bm{x}}_m^\top$ (default: \emph{no}
centring, see Remark~\ref{rem:center}). Take the (thin) SVD
$\tilde{X}_m = U_m \Sigma_m V_m^\top$ and define the \emph{window frame}
\begin{equation}
F_m \;=\; \big[\,\bm{v}_1 \; \bm{v}_2 \; \cdots \; \bm{v}_k\,\big]\in\R^{N\times k},
\qquad V_m = [\bm{v}_1,\bm{v}_2,\dots],
\end{equation}
i.e.\ the top-$k$ right singular vectors. Its columns are orthonormal and span the
$k$-dimensional subspace of neuron space that captures maximal window variance;
$\Span(F_m)$ is a point of $\Gr(k,N)$. Each frame is time-stamped at the window
centre $c_m = a_m + \lfloor w/2\rfloor$.
\end{definition}

\begin{definition}[Reliability diagnostics]
For window $m$ let $\{\sigma_i\}=\{\sigma_i(\tilde{X}_m)\}$. The \emph{explained-variance
ratio} and the \emph{singular-value gap} are
\begin{equation}
\mathrm{evr}_m = \frac{\sum_{i\le k}\sigma_i^2}{\sum_i \sigma_i^2},
\qquad
g_m = \frac{\sigma_k}{\sigma_{k+1}}.
\end{equation}
A frame is well-defined only where $g_m \gg 1$; as $g_m\to 1$ the $k$-th direction is
indistinguishable from the $(k{+}1)$-th and the point on $\Gr(k,N)$ is dominated by
noise. We report $g_m$ alongside every kinematic curve.
\end{definition}

\begin{remark}[Uncentred vs.\ centred windows]
\label{rem:center}
Uncentred SVD extracts the subspace the activity \emph{occupies} (it includes the
window mean direction), which is stable even at rest. Centred SVD extracts the local
\emph{tangent/velocity} subspace, which is noise-dominated when the state is nearly
stationary and does not discriminate conditions. The default is uncentred.
\end{remark}

The output of this step, per (trial, $k$), is a frame sequence
$\{F_m\}_{m=1}^{M}$, times $\{t_{c_m}\}$, diagnostics $\{g_m\}$, and the geodesic
distance matrix of \S\ref{sec:geom}.

\section{Grassmannian geometry}
\label{sec:geom}

\subsection{Manifold and representations}
The Grassmannian $\Gr(k,N)$ is the set of $k$-dimensional linear subspaces of
$\R^N$. Two representations are used:
\begin{enumerate}[leftmargin=1.4em]
\item \textbf{Frame:} an orthonormal $F\in\R^{N\times k}$ with $\Span(F)$ the point.
Frames are non-unique: $F$ and $FQ$ ($Q\in O(k)$) represent the same subspace, so
$\Gr(k,N) = \mathrm{St}(k,N)/O(k)$ (Stiefel quotient).
\item \textbf{Projector:} the orthogonal projector $P = FF^\top\in\R^{N\times N}$.
This is unique: $P=P^\top$, $P^2=P$, $\tr P = k$. Geometry is computed on projectors,
so quotient ambiguity never appears.
\end{enumerate}
The manifold dimension is $\dim\Gr(k,N)=k(N-k)$. For $k=1$, a point is a line through
the origin, and $\Gr(1,N)=\RP^{N-1}$ (real projective space); this identification
underlies the topological signatures of \S\ref{sec:tda}.

\subsection{Principal angles and distances}
\begin{definition}[Principal angles]
For frames $A,B\in\R^{N\times k}$ let $\sigma_1\ge\cdots\ge\sigma_k$ be the singular
values of $A^\top B$ (all in $[0,1]$). The \emph{principal angles} are
$\theta_i = \arccos\sigma_i \in [0,\tfrac{\pi}{2}]$, $i=1,\dots,k$.
\end{definition}

The principal angles are the complete $O(k)$-invariant of a pair of subspaces and
generate every metric we use:
\begin{align}
d_{\mathrm{arc}}(A,B) &= \Big(\textstyle\sum_{i} \theta_i^2\Big)^{1/2}
&&\text{(arc-length / geodesic),}\label{eq:darc}\\
d_{\mathrm{can}}(A,B) &= \sqrt{2}\,\, d_{\mathrm{arc}}(A,B)
&&\text{(canonical, \texttt{geomstats}),}\label{eq:dcan}\\
d_{\mathrm{ch}}(A,B) &= \Big(\textstyle\sum_i \sin^2\theta_i\Big)^{1/2}
= \tfrac{1}{\sqrt 2}\,\norm{P_A - P_B}_F
&&\text{(chordal / projection).}\label{eq:dch}
\end{align}
Equations \eqref{eq:darc}--\eqref{eq:dcan} differ only by the global factor
$\sqrt 2$; because this is a uniform rescaling of all distances, \emph{persistent
homology, Betti numbers, distance-ratio rankings, and curve \emph{shapes} are
invariant} to the choice, while absolute magnitudes are not. The chordal distance
\eqref{eq:dch} is the Euclidean distance between projectors (up to $\sqrt2$) and
agrees with the geodesic distance to second order in the angles,
$\sin\theta = \theta - \theta^3/6 + O(\theta^5)$, but saturates as
$\theta\to\pi/2$. We use $d_{\mathrm{can}}$ as the default distance for
kinematics/TDA and report the chordal--geodesic relation as a diagnostic
(\S\ref{sec:tpca}). The pairwise geodesic distance matrix of a frame sequence is
$D_{mn} = d_{\mathrm{can}}(F_m,F_n)$.

\begin{remark}[Fast path]
Because \eqref{eq:darc} needs only the $k\times k$ SVD of $A^\top B$, the distance
matrix costs $O(M^2 k^2 N)$, avoiding $N\times N$ projector operations. This fast
path is validated to equal the \texttt{geomstats} canonical metric to $\sim10^{-15}$.
\end{remark}

\subsection{Tangent space, geodesics, transport}
At a point $P$ (projector form) the tangent space consists of symmetric matrices
that infinitesimally rotate $\Span(P)$ into its orthogonal complement (the
``horizontal'' directions). The Riemannian \emph{logarithm}
$\Log_P(Q)\in T_P\Gr(k,N)$ is the initial velocity of the unit-time geodesic from $P$
to $Q$, and the \emph{exponential} $\Exp_P(V)$ is its inverse; they satisfy
$\norm{\Log_P(Q)}_P = d_{\mathrm{can}}(P,Q)$, where $\norm{\cdot}_P$ is the norm
induced by the canonical metric $g_P(\cdot,\cdot)$. Parallel transport
$\PT_{P\to Q}:T_P\to T_Q$ moves a tangent vector along the connecting geodesic while
preserving inner products. We use the closed-form \texttt{geomstats} implementations
of $\Log$, $\Exp$, $g_P$, and $\PT$.

\section{Riemannian kinematics}
\label{sec:kin}

Treat the frame sequence as a discrete curve $P_1,\dots,P_M$ on $\Gr(k,N)$ with
uniform time step $\Delta t = t_{c_{m+1}}-t_{c_m}$.

\begin{definition}[Velocity and speed]
\begin{equation}
\bm{v}_m = \frac{1}{\Delta t}\,\Log_{P_m}(P_{m+1}) \in T_{P_m}\Gr(k,N),
\qquad
s_m = \norm{\bm{v}_m}_{P_m} = \frac{d_{\mathrm{can}}(P_m,P_{m+1})}{\Delta t}.
\end{equation}
The speed is thus the geodesic distance travelled per unit time; it measures how fast
the occupied subspace reorients.
\end{definition}

\begin{definition}[Covariant acceleration]
Because $\bm{v}_{m-1}$ and $\bm{v}_m$ live in \emph{different} tangent spaces, they
must be compared after transport:
\begin{equation}
\bm{a}_m = \frac{1}{\Delta t}\Big(\bm{v}_m - \PT_{P_{m-1}\to P_m}\bm{v}_{m-1}\Big)
\in T_{P_m}\Gr(k,N),
\qquad
\alpha_m = \norm{\bm{a}_m}_{P_m}.
\end{equation}
This is the discrete covariant derivative $\tfrac{D\bm{v}}{dt}$. A geodesic has
$\bm{a}\equiv 0$; hence $\alpha_m$ quantifies how far the subspace path deviates from
``coasting in a straight line'' on the manifold. (A naive difference
$\bm{v}_m-\bm{v}_{m-1}$ without transport mixes tangent spaces and is not
intrinsically meaningful.)
\end{definition}

\begin{definition}[Geometric curvature]
Split the acceleration into components parallel and orthogonal to the velocity,
$\bm{a}_m = \bm{a}_m^{\parallel} + \bm{a}_m^{\perp}$ with
$\bm{a}_m^{\parallel} = g_{P_m}(\bm{a}_m,\hat{\bm{v}}_m)\,\hat{\bm{v}}_m$,
$\hat{\bm{v}}_m = \bm{v}_m/s_m$. The path curvature is
\begin{equation}
\kappa_m = \frac{\norm{\bm{a}_m^{\perp}}_{P_m}}{s_m^{2}} .
\end{equation}
$\kappa_m$ is reparametrisation-invariant (it depends on the trace of the path, not
its speed) and equals the inverse radius of the osculating circle.
\end{definition}

\begin{definition}[Global summaries]
Path length, endpoint displacement and geodesic efficiency are
\begin{equation}
L = \sum_{m} d_{\mathrm{can}}(P_m,P_{m+1}),\qquad
\Delta = d_{\mathrm{can}}(P_1,P_M),\qquad
\mathcal{E} = \frac{\Delta}{L}\in(0,1].
\end{equation}
$\mathcal{E}\to 1$ indicates a near-geodesic, directed reconfiguration; $\mathcal{E}\to 0$
indicates a closed or space-filling orbit (a returning loop has $\Delta\approx 0$ but
$L$ large). Constant $s_m$, constant $\kappa_m$, and constant nonzero $\alpha_m$
together certify uniform circular motion on the manifold.
\end{definition}

\section{Intrinsic statistics of the subspace distribution}
\label{sec:tpca}

\begin{definition}[Fr\'echet (Karcher) mean]
The intrinsic mean of frames $\{P_i\}$ is
$\bar P = \arg\min_{P} \sum_i d_{\mathrm{can}}(P,P_i)^2$, computed by the fixed-point
iteration
$P \leftarrow \Exp_P\!\big(\tfrac{1}{n}\sum_i \Log_P(P_i)\big)$ until the mean
tangent step falls below tolerance.
\end{definition}

\begin{definition}[Tangent-space PCA]
Map every frame to the tangent space at $\bar P$ and vectorise,
$\bm{\ell}_i = \operatorname{vec}\big(\Log_{\bar P}(P_i)\big)\in\R^{N^2}$. Centre and
take the SVD $L-\mathbf{1}\bar{\bm{\ell}}^\top = U\Sigma W^\top$; the principal
directions $W_{\cdot j}$ are \emph{modes of subspace variation}, the scores
$U\Sigma$ are coordinates, and the explained-variance ratios
$\lambda_j = \sigma_j^2/\sum_{j'}\sigma_{j'}^2$ define an \emph{intrinsic
dimensionality}: the number of tangent PCs needed to reach, say, $90\%$ cumulative
variance. This is the manifold analogue of PCA applied to the subspace trajectory
itself (PCA of subspaces, not of states).
\end{definition}

The chordal--geodesic scatter $\{(d_{\mathrm{arc}},d_{\mathrm{ch}})\}$ over sampled
frame pairs visualises the small-angle regime (where the two agree) versus the
large-angle regime (where chordal saturates below $y=x$); systematic bending below
the diagonal certifies that the sampled subspaces span a wide angular range rather
than clustering.

\section{Discrete exterior calculus on the sampled submanifold}
\label{sec:dec}

To equip the sampled subspace ensemble with a differential structure we build a
$2$-manifold surrogate and compute intrinsic operators on scalar fields.

\paragraph{Complex construction.} Pool frames across all trials of a condition;
compute the geodesic distance matrix $D$; embed to $\R^2$ by metric multidimensional
scaling (MDS); Delaunay-triangulate the planar points to obtain a simplicial
$2$-manifold $\mathcal{K}$ (vertices $V$, edges $E$, triangles $T$). A long-edge
prune ($>\!p$-th percentile) yields an optional ``hole-hint'' complex $\mathcal{K}'$
(a poor-man's alpha complex); the DEC operators, however, run on the full Delaunay
disk, which is a clean valid $2$-manifold.

\paragraph{Cochains and operators.} A scalar field is a $0$-cochain $f\in C^0$ with
one value per vertex. Discrete exterior calculus provides
\begin{align}
\text{exterior derivative} \quad & d: C^p\to C^{p+1}, &&(df)(\text{edge }[ij]) = f_j - f_i,\\
\text{Hodge star} \quad & \star: C^p\to C^{n-p}, && \text{(mass/volume weighting)},\\
\text{codifferential} \quad & \delta = \pm\,\star d\,\star, &&\\
\text{Laplace--de Rham} \quad & \Delta = d\delta + \delta d, && \Delta f = \delta d f \ \ (p{=}0).
\end{align}
Here $df$ is the discrete gradient $1$-form (a value per edge); $\norm{df}$ localises
steep gradients of $f$ over the subspace manifold; $\Delta f$ localises sources and
sinks (concentrations) of $f$. The sharp operator $\sharp$ turns $df$ into a vertex
vector field for streamlines. Scalar fields carried on the vertices are
\begin{equation}
\underbrace{\tfrac1N\sum_i x_i^2}_{\text{energy}},\quad
\underbrace{s_m}_{\text{subspace speed}},\quad
\underbrace{\frac{(\sum_i\lambda_i)^2}{\sum_i\lambda_i^2}}_{\text{participation ratio}},\quad
\underbrace{\text{input drive}}_{\text{stimulus}},\quad
\underbrace{\frac{\norm{F^\top\bm{x}}^2}{\norm{\bm{x}}^2}}_{\text{self-capture}} .
\end{equation}
The \emph{self-capture} of a state $\bm{x}$ by its frame $F$ equals the fraction of
state energy inside the subspace, $\cos^2\!\angle(\bm{x},\Span F)$ for $k=1$, i.e.\
$1-d_{\mathrm{ch}}^2$ between the state line and $\Span F$.

\paragraph{Betti numbers (self-contained).} From the oriented boundary matrices
$\partial_1\in\R^{V\times E}$ (edges$\to$vertices) and
$\partial_2\in\R^{E\times T}$ (triangles$\to$edges),
\begin{equation}
b_0 = V - \rank\partial_1,\qquad
b_1 = E - \rank\partial_1 - \rank\partial_2,\qquad
b_2 = T - \rank\partial_2,
\end{equation}
giving the number of connected components, independent loops, and enclosed voids of
$\mathcal{K}$ (resp.\ $\mathcal{K}'$).

\paragraph{Cross-projection matrix.} For one trial, with states $\bm{x}(t_i)$ at
window centres and all visited frames $\{F_j\}$,
\begin{equation}
R_{ij} = \frac{\norm{F_j^\top \bm{x}(t_i)}^2}{\norm{\bm{x}(t_i)}^2} \in [0,1],
\end{equation}
the fraction of the state at time $t_i$ captured by the subspace visited at time
$t_j$. Its diagonal is the self-capture time course; its off-diagonal stripes image
recurrence of the subspace (a loop produces periodic bands parallel to the diagonal).

\begin{remark}[Read DEC qualitatively]
The MDS embedding is not isometric, so $\Delta f$ magnitudes are inflated and
should be read qualitatively (sign/pattern), while $df$ and the Betti numbers are
robust. Rigorous loop topology is deferred to persistent homology
(\S\ref{sec:tda}), which needs no embedding.
\end{remark}

\section{Persistent homology}
\label{sec:tda}

Persistent homology is applied directly to distance matrices, so it never requires a
Euclidean embedding and is insensitive to the global $\sqrt2$ scale.

\begin{definition}[Filtration and diagrams]
Given a finite metric space with distance matrix $D$, the Vietoris--Rips complex at
scale $\varepsilon$ has a simplex on any subset with pairwise distances
$\le\varepsilon$. Growing $\varepsilon$ from $0$ produces a nested filtration; each
homology class (component in $H_0$, loop in $H_1$, void in $H_2$) is born at some
$\varepsilon_b$ and dies at $\varepsilon_d$. The multiset of points
$(\varepsilon_b,\varepsilon_d)$ is the \emph{persistence diagram}; the
\emph{lifetime} $\varepsilon_d-\varepsilon_b$ measures feature robustness.
\end{definition}

We compute $H_0,H_1,H_2$ over the field $\Fld_2$ (coefficients mod $2$). $\Fld_2$ is
chosen deliberately: it is the field under which \emph{non-orientable / projective}
structure is visible, and $\Gr(1,N)=\RP^{N-1}$ is non-orientable, so a rotating line
direction registers as an $H_1$ class over $\Fld_2$. Two scopes are analysed:
\begin{itemize}[leftmargin=1.4em]
\item \textbf{single}: the within-trial trajectory distance matrix $D$ (detects
loops traced \emph{within} one trial, e.g.\ a moving-stimulus orbit);
\item \textbf{pooled}: the across-trial ensemble cloud (detects structure that only
exists \emph{across} conditions, e.g.\ a settling ring sampled trial-by-trial).
\end{itemize}

\begin{definition}[Bottleneck distance]
For two diagrams $\mathcal{D}_1,\mathcal{D}_2$ (in a fixed dimension, here $H_1$), the
bottleneck distance is
$d_B = \inf_{\gamma}\sup_{p} \norm{p-\gamma(p)}_\infty$ over partial matchings
$\gamma$ (unmatched points sent to the diagonal). We report the pairwise $H_1$
bottleneck matrix across conditions as a topological ``fingerprint'' dissimilarity.
\end{definition}

\section{Direct-manifold cross-check}
\label{sec:direct}

To distinguish genuine structure from artefacts of the subspace map, the identical
persistent homology is run on the raw state manifold. Per trial we form the Euclidean
state distance matrix (optionally subsampled), embed by PCA$(3)$ and Isomap$(2)$ for
visualisation, and compute $H_0/H_1/H_2$ over $\Fld_2$. Because the extrinsic
Euclidean state distance and the intrinsic Grassmannian geodesic distance are
different objects, agreement of the two $H_1$ signatures --- quantified by the
bottleneck distance after scaling each diagram to unit maximum lifetime --- is strong
evidence that a detected loop is real rather than a metric artefact.

\section{Load-bearing conventions}
\label{sec:conv}

\begin{enumerate}[leftmargin=1.4em]
\item \textbf{Metric scale.} $d_{\mathrm{can}}=\sqrt2\,d_{\mathrm{arc}}$; topology and
curve shapes are scale-invariant, absolute magnitudes are not.
\item \textbf{$k$ as a topological filter.} Different $k$ expose different structure.
Example: a rotating line direction produces an $H_1$ loop at $k=1$ (on $\RP^{N-1}$)
but not at $k=2$ if the containing $2$-plane is static. Always analyse multiple $k$.
\item \textbf{Frame reliability.} A $k$-frame is trustworthy only where
$\sigma_k/\sigma_{k+1}\gg1$. Locally rank-$r$ dynamics make $k>r$ fit noise (huge
subspace speed, spurious homology).
\item \textbf{Windows uncentred} by default (occupied subspace, not tangent
subspace).
\item \textbf{Regime placement.} Loops may live within a single trial (moving
stimulus) or only across trials (settling attractor); analyse both scopes.
\end{enumerate}

\section*{Estimator summary}
\addcontentsline{toc}{section}{Estimator summary}
\begin{center}
\small
\begin{tabular}{@{}lll@{}}
\toprule
Quantity & Symbol / formula & Interpretation \\
\midrule
Window frame & top-$k$ right sing.\ vectors of $\tilde X_m$ & occupied subspace \\
Reliability & $g_m=\sigma_k/\sigma_{k+1}$ & frame trustworthiness \\
Geodesic distance & $d_{\mathrm{can}}=\sqrt2(\sum\theta_i^2)^{1/2}$ & subspace dissimilarity \\
Speed & $s_m=d_{\mathrm{can}}(P_m,P_{m+1})/\Delta t$ & reorientation rate \\
Covariant accel. & $\alpha_m=\norm{(\bm v_m-\PT\bm v_{m-1})/\Delta t}$ & deviation from geodesic \\
Curvature & $\kappa_m=\norm{\bm a^\perp}/s_m^2$ & path bending \\
Efficiency & $\mathcal E=\Delta/L$ & directed vs.\ closed orbit \\
Intrinsic dim. & \#tangent-PCs for $90\%$ var. & subspace-variation complexity \\
Self-capture & $\norm{F^\top\bm x}^2/\norm{\bm x}^2$ & state energy in subspace \\
Betti & $b_0,b_1,b_2$ & components / loops / voids \\
Persistence & lifetime $\varepsilon_d-\varepsilon_b$ ($\Fld_2$) & robust topology \\
Fingerprint & $H_1$ bottleneck $d_B$ & condition dissimilarity \\
\bottomrule
\end{tabular}
\end{center}

\end{document}
